\chapter{The architecture}
\label{ch:arch}

\section{Nine programs, not one}

When you run IncidentFlow, nine separate programs start. That number surprises
people. A student project is usually one program; this is nine, and every one of
them earns its place.

Here is the whole system on one page. Read the arrows as ``talks to''.

\begin{figure}[H]
\centering
\begin{tikzpicture}[
  font=\footnotesize\sffamily,
  box/.style={draw=rule,fill=white,rounded corners=3pt,minimum height=9mm,minimum width=24mm,align=center},
  edge/.style={draw=rule,fill=brandlt,rounded corners=3pt,minimum height=9mm,minimum width=24mm,align=center},
  data/.style={draw=rule,fill=goodbg,rounded corners=3pt,minimum height=9mm,minimum width=24mm,align=center},
  work/.style={draw=rule,fill=warnbg,rounded corners=3pt,minimum height=9mm,minimum width=24mm,align=center},
  arr/.style={-{Stealth[length=2mm]},draw=muted,thick},
  evt/.style={-{Stealth[length=2mm]},draw=muted,thick,dashed},
]
\node[box] (browser) {Your browser};
\node[edge, below=9mm of browser] (nginx) {\textbf{nginx}\\port 8080};
\node[box, below left=11mm and 6mm of nginx] (web) {\textbf{web}\\React files};
\node[box, below=11mm of nginx] (api) {\textbf{api}\\Laravel + PHP-FPM};
\node[box, below right=11mm and 6mm of nginx] (rt) {\textbf{realtime}\\Node, SSE/WS};

\node[data, below=11mm of api] (pg) {\textbf{postgres}\\the truth};
\node[data, left=8mm of pg] (redis) {\textbf{redis}\\queue + events};
\node[work, right=8mm of pg] (horizon) {\textbf{horizon}\\job worker};
\node[work, below=7mm of redis] (sched) {\textbf{scheduler}\\timed tasks};
\node[box, below=7mm of horizon] (mail) {\textbf{mailpit}\\fake inbox};

\draw[arr] (browser) -- (nginx);
\draw[arr] (nginx) -- (web);
\draw[arr] (nginx) -- (api);
\draw[arr] (nginx) -- (rt);
\draw[arr] (api) -- (pg);
\draw[arr] (api) -- (redis);
\draw[evt] (redis) -- (rt);
\draw[arr] (horizon) -- (redis);
\draw[arr] (horizon) -- (pg);
\draw[arr] (sched) -- (redis);
\draw[arr] (horizon) -- (mail);
\end{tikzpicture}
\caption*{\small\color{muted}The nine containers. Solid arrows are direct
connections; the dashed arrow is events flowing from Laravel to the realtime
tier through Redis. Only nginx is reachable from outside.}
\end{figure}

\section{What each one is for}

\begin{center}
\begin{tabular}{@{}>{\raggedright\arraybackslash}p{2.4cm}>{\raggedright\arraybackslash}p{2.5cm}>{\raggedright\arraybackslash}p{9.0cm}@{}}
\toprule
\textbf{Container} & \textbf{Built from} & \textbf{Its job} \\
\midrule
\fpath{nginx} & nginx 1.27 & The front door. Everything from your browser arrives here and is routed onward. Nothing else is exposed. \\
\fpath{web} & React + Vite & Serves the compiled JavaScript and CSS. It has no logic --- it asks the API for everything. \\
\fpath{api} & PHP 8.3 + Laravel & All the rules. Who may do what, what a valid transition is, what gets written. \\
\fpath{realtime} & Node.js 22 & Holds long-lived connections open so screens update without a refresh. \\
\fpath{postgres} & PostgreSQL 16 & Stores everything. The single source of truth. \\
\fpath{redis} & Redis 7 & Two jobs: carries live events from Laravel to realtime, and holds the queue of pending work. \\
\fpath{horizon} & same as api & Runs queued jobs --- chiefly sending email --- so the web request does not wait for them. \\
\fpath{scheduler} & same as api & Runs timed tasks, like sweeping notifications that were never delivered. \\
\fpath{mailpit} & mailpit & A fake inbox for development, so test email never reaches a real person. \\
\bottomrule
\end{tabular}
\end{center}

\note{Notice that \fpath{api}, \fpath{horizon} and \fpath{scheduler} are built
from \emph{the same image}. One Dockerfile, three containers, three different
start commands. Without Docker you would install PHP once and run three
processes; here each gets its own isolated process, so a memory leak in the
queue worker cannot take the API down with it.}

\section{Following one request all the way through}

Abstract diagrams are easy to nod at and hard to remember. So follow one real
action: \textbf{a responder clicks ``Acknowledge'' on an incident.}

\begin{enumerate}[leftmargin=1.5cm]
\item The browser sends \fpath{POST /api/v1/incidents/42/status} to
  \fpath{localhost:8080}. That is \textbf{nginx}.
\item nginx sees the path starts with \fpath{/api/}, so it hands the request to
  \textbf{api} over FastCGI --- not HTTP. (FastCGI is the protocol web servers
  use to talk to PHP.)
\item Laravel checks the JWT in the \fpath{Authorization} header: who is this?
\item It checks the \emph{policy}: may a responder change this incident's
  status?
\item It checks the \emph{enum}: is \fpath{open} to \fpath{acknowledged} a legal
  move?
\item Inside one database transaction it writes the new status to
  \textbf{postgres}, appends a timeline event, and creates notification rows.
\item The transaction commits. Only then does it push jobs onto \textbf{redis}
  and publish an event.
\item \textbf{horizon} picks the job off Redis and sends email through
  \textbf{mailpit}.
\item \textbf{realtime} is subscribed to Redis. It receives the event and pushes
  it down every open connection.
\item Every other open browser updates. Nobody pressed refresh.
\end{enumerate}

\warnbox{Step 7 says ``the transaction commits, \emph{only then}'' for a
reason. If the job were queued before the commit, Horizon --- a separate process
--- could pick it up and look for an incident that has not been written yet. It
would find nothing and fail. This is one of the classic distributed-system bugs,
and the ordering in that step is the entire defence against it.}

\section{Why not one program?}

You could build all of this as a single Laravel application. Many people would.
Three things here genuinely cannot be done that way.

\begin{description}[leftmargin=1.2cm,style=nextline]
\item[Long-lived connections] PHP-FPM assigns a worker process per request. A
  connection held open for an hour would occupy a worker for an hour. Twenty
  users watching a dashboard would consume the whole pool and the API would stop
  answering anything at all. Node handles thousands of idle connections on one
  process, because that is what an event loop is good at. Hence a separate
  \fpath{realtime} service.
\item[Slow work] Sending an email takes as long as the mail server takes. If
  that happened inside the request, the responder's browser would spin while
  SMTP negotiated. So the API writes a row, returns immediately, and
  \fpath{horizon} does the sending afterwards.
\item[Timed work] Something must run every minute regardless of whether anyone
  is using the system. A web request cannot do that, because there might not be
  one. Hence \fpath{scheduler}.
\end{description}

\note{This is a \emph{modular service-oriented} architecture, not microservices.
The distinction matters: there is one database and one deployable unit. The
services are split along the lines where the \emph{runtime model} differs ---
request/response, long-lived connection, background work --- not along business
domains. Splitting by domain is what produces microservices, and with them
distributed transactions, which this system deliberately does not have.}

\section{The security boundary}

One rule, and it explains most of the configuration you will read later:
\textbf{only nginx is reachable from outside}.

\begin{center}
\begin{tabular}{@{}lll@{}}
\toprule
\textbf{Service} & \textbf{Reachable from your machine?} & \textbf{How} \\
\midrule
nginx & yes & \fpath{localhost:8080} \\
mailpit & yes (development only) & \fpath{localhost:8025} \\
postgres & development only & \fpath{localhost:5432} \\
redis & development only & \fpath{localhost:6379} \\
api & no & only nginx, over the internal network \\
realtime & no & only nginx \\
web & no & only nginx \\
horizon & no & has no network port at all \\
scheduler & no & has no network port at all \\
\bottomrule
\end{tabular}
\end{center}

In production those ``development only'' rows are closed too. You will see
exactly how in Part~III, in \fpath{docker-compose.prod.yml}.
